\subsection{Arah Penyelesaian Masalah}
\label{subsec:arah-penyelesaian-masalah}

Berdasarkan masalah yang telah diidentifikasi pada Subbab \ref{subsec:identifikasi-masalah}, solusi yang diperlukan bukan hanya penggantian format token, tetapi pembaruan cara DecMed mendefinisikan, merepresentasikan, menegakkan, dan mendelegasikan hak akses. DecMed sudah memiliki fondasi yang sesuai untuk sistem RME terdesentralisasi melalui IOTA, IPFS, dan PRE. Oleh karena itu, solusi yang dianalisis pada bagian ini diarahkan untuk memperluas lapisan kontrol akses tanpa mengubah peran utama komponen-komponen tersebut.

Berdasarkan kebutuhan sistem pada Subbab \ref{subsec:kebutuhan-sistem}, solusi yang dianalisis dikelompokkan ke dalam aspek utama berikut: model kebijakan akses, mekanisme token otorisasi, segmentasi objek data, delegasi akses, penegakan otoritas, validasi kewenangan aktor, serta revokasi dan audit delegasi. Pengelompokan ini diperlukan karena kontrol akses yang lebih luas tidak dapat ditegakkan hanya melalui satu komponen. Model kebijakan diperlukan untuk mendefinisikan hubungan antara aktor, objek, hak akses, dan kondisi; token otorisasi diperlukan untuk membawa kapabilitas; segmentasi data diperlukan agar kebijakan memiliki objek penegakan yang jelas; sedangkan \textit{Server} PRE diperlukan sebagai titik penegakan sebelum data terenkripsi diambil dan dire-enkripsi.

Aspek-aspek tersebut saling bergantung. Kebijakan yang lebih luas tidak cukup apabila data yang dilindungi masih berbentuk objek terenkripsi yang terlalu besar. Sebaliknya, segmentasi data tidak dapat dimanfaatkan secara aman apabila kebijakan akses belum mendefinisikan siapa yang boleh mengakses setiap objek dan dalam kondisi apa. Delegasi akses juga memerlukan media kapabilitas yang dapat dipersempit tanpa membuka peluang perluasan hak akses, sedangkan revokasi dan audit diperlukan agar delegasi yang terjadi tetap dapat dikendalikan dan ditelusuri. Dengan demikian, analisis solusi pada bagian ini berfungsi sebagai dasar pemilihan pendekatan sebelum rancangan teknis dijelaskan pada Subbab \ref{sec:rancangan-solusi}.
